\section{Open Source 的严格含义：从“可用”到“可复现”}

\subsection{定义争议的技术本质}
讲者后半段讨论了开放定义的边界问题：如果只有权重、缺少训练数据与完整流程，是否足以称为 open source？这不是语义争执，而会影响科研可复现性和责任划分。

\begin{importantbox}{Open Source for Models 需要新的定义框架}
传统软件开源主要围绕“源代码可见可改”。  
模型系统中，数据、训练流程、评测协议与发布策略同样是“源”的组成部分。
\end{importantbox}

\subsection{可复现性阶梯}
\begin{table}[H]
\centering
\begin{tabular}{p{3.0cm}p{3.8cm}p{4.8cm}}
\toprule
\textbf{开放层级} & \textbf{可做的复现} & \textbf{仍然缺失} \\
\midrule
API 可用 & 行为复测、接口级评测 & 机制归因、训练追溯 \\
Weights 可得 & 推理复现、微调复现、部分解释实验 & 数据与训练因果链 \\
Full source & 训练与评测端到端复现、系统级改造 & 主要瓶颈转为算力与工程资源 \\
\bottomrule
\end{tabular}
\caption{从“可调用”到“可复现”的能力阶梯}
\end{table}

\begin{knowledgebox}{为什么“可复现”比“可下载”更关键}
科学共同体需要可验证的知识积累。只有当关键实验条件可复现、可比对、可反驳，领域才能形成稳固共识，而非靠个别组织的不可验证声明。
\end{knowledgebox}

\begin{warningbox}{把“开源”当作营销标签会反噬社区信任}
若术语长期被滥用，研究者会在协作预期上出现系统性误判：  
以为可复现，实际不可复现；以为可审计，实际不可审计。
\end{warningbox}

\subsection{本章小结}
本章核心是“定义即基础设施”。对于模型系统，开放定义直接决定研究可复现性、协作效率与治理可操作性。


